% *- coding: UTF-8 -*-
% !TEX builder = Texlive
% !TEX prggram = xelatex
\documentclass[UTF8,a4paper,12pt,titlepage]{ctexart}
% tittlepage制作一个封面
% draft是草稿，有个方框
%\documentclass[UTF8,a4paper,11pt]{book}

%\CTEXsetup[name={第,章}]{section}%每section统一名称

%每章节题目统一格式
\CTEXsetup[format={\zihao{-3}\raggedright\bfseries}]
{section}

\newcommand{\keyword}[1]{\textbf{#1}}

%\usepackage{microtype}

% 纸张格式
\usepackage[paper=a4paper,inner=1.5cm,outer=2cm,top=2cm,
bottom=2.5cm,bindingoffset=0.5cm]{geometry}

\usepackage[onehalfspacing]{setspace} % 1.5倍行距

% 自动链接内容
\usepackage[colorlinks=true,linkcolor=black]{hyperref}
% 设置文件属性内容
\hypersetup{pdfauthor={吴毓明},
            pdftitle={储罐底部静载称重模块的深化设计},
            pdfsubject={使用说明},
            pdfkeywords={任何问题，联系61508712@qq.com}}

% 定义页眉页脚
%\include{foot_head}

% 表格 定义表头字体
\newcommand{\head}[1]{\textnormal{\textbf{#1}}}

% 表格 特殊命令，见后文 array
% 在竖向多种对齐方式
\usepackage{array}
\setlength{\extrarowheight}{4pt} % 增加列距
\usepackage{booktabs} % 用于美化表格的线
\usepackage{multirow} % 合并行
\newcommand{\normal}[1]{\multicolumn{1}{l}{#1}}

% 图片
\usepackage{graphicx}
\usepackage{wrapfig} % 文字环绕图片

% 流程图
% 导入tikz绘图包
\usepackage{tikz}
% 导入基本图形
\usetikzlibrary{shapes,arrows}
%说明文字
\tikzstyle{text_box} = [
   rectangle,
   rounded corners,
   minimum width=1.5cm,
   minimum height=0.5cm,
   text centered, 
   ]
%箭头
\tikzstyle{arrow} = [thick,->,>=stealth]

%多行公式
\usepackage{amsmath}

\begin{document}

\title{
   \zihao{1} 高锰酸钾 投加系统 加药量 计算\\[2.5mm]
   {\noindent}\rule{16cm}{1pt}\\[3mm]
   % \zihao{2} 广州市利亨昌贸易有限公司\\
   \zihao{2} 宁夏宁东水务有限责任公司\\
   % \zihao{4} 吴毓明\\
   %\zihao{4} Wym\\
   %\zihao{4} 2022年12月27日
   %\begin{figure}[h]
   %   \centering
   %   \includegraphics[height=10cm]{g05.jpg}
   %\end{figure}
}

\maketitle % 封面

%\tableofcontents % 目录

\newpage % 换页

%\part{系统组成}
%ctex的\chapter{sdsg}是不能用的

\section{设计要求}
    本项目设计 使用 高锰酸钾制备投加一体化装置 制备成 3\% 的 高锰酸钾 溶液 投加，
    最大 投加量 按 1mg/L计，
    原液浓度按 32\%计，
    调配次数近期按每天2次计，
    远期按每天3次计；
    近期浓碱输液泵2台，
    2用1备；高锰酸钾溶液 投药泵3台，

    投加位置及投加量如下：

      \begin{enumerate}
          \item 进水泵房：2 处，1mg/L（每处规模为 40 万 $m^3/d$）；
      \end{enumerate}


\section{投加量计算}
    \subsubsection{计算推导}
        \par根据 《化学化工 物性数据手册》（化学工业出版社）无机卷，
        第555页，
        3\%（质量分数）的 高锰酸钾 水溶液 在 15 $^{\circ}C$ 时，
        密度是 $\rho_{3\%KMnO_4}=1.0200 g/cm^3=1.0200kg/L$。
      \begin{figure}[h]
         \centering
         \includegraphics[height=3.5cm]{p1.PNG}
         \caption{3\% 高锰酸钾 溶液物性数据}
      \end{figure}\\

        设密度为$\rho$、
        质量为$m$、
        体积为$V$、
        质量流量为$Q_m$、
        体积流量为$Q_V$、
        时间为$T$，则：
        \begin{align}
            \rho &= m \div V \\
            Q_m &= m \div T\\
            Q_v &= V \div T
        \end{align}
        由$\rho = m \div v$可得$m = \rho * V$，
        代入$Q_m = m \div T$可得：
        \begin{equation}
            Q_m = \rho * V \div T = \rho * Q_v
        \end{equation}
        则：
        \begin{equation}
            Q_v = Q_m \div \rho = Q_m \div \rho_{3\%KMnO_4}
        \end{equation}\label{equation5}

    \subsubsection{进水泵房 投加量计算}
        设计水量为40万$m^3/d$，
        高锰酸钾 的 投率 为1mg/L，
        则 高锰酸钾 溶质 的 投加量 为：\\
        \begin{equation}
            400000m^3/d \times 1mg/L = \frac{400000}{24} m^3/h
            \times \frac{1 \times 0.001 \times 0.001}{0.001}kg/m^3
            = 16.6666666667 kg/h
        \end{equation}
        投加 高锰酸钾 溶液的 浓度（质量分数）为3\%，
        则 投加的 质量流量 $Q_m$为：
        \begin{equation}
            Q_m = 16.6666666667 kg/h \div 3\% = 555.5555555556 kg/h
        \end{equation}
        代入\ref{equation5}，可得：
        \begin{equation}
            Q_v = Q_m \div \rho_{3\%KMnO_4} = 555.5555555556 kg/h \div 1.0200g/cm^3 = 544.6623093682 L/h
        \end{equation}
        厂区进水总管 3\%高锰酸钾溶液 设计投加量 约为 544.6623093682L/h。
        \par此投加量由2台投加泵共同投加，
        则每台投加泵投加量为：
        \begin{equation}
            544.6623093682 L/h \div 2 = 272.3311546841 L/h
        \end{equation}


\end{document}
